\section{Introduction}\label{sec:intro}
Space-cooling electricity demand in India is growing rapidly as room air-conditioner
ownership rises and urban heat intensifies \cite{iea2023}. Pune's composite climate
places three different demands on residential cooling: a hot, dry pre-monsoon with
afternoon peaks near 41\,$^\circ$C, a humid monsoon in which latent loads and
ventilation dominate, and a mild winter that needs no cooling at all. Most homes are
still controlled by the on/off thermostat of a split air conditioner. That thermostat
ignores three things: time-of-day (ToD) tariffs; the thermal inertia of reinforced-concrete
roofs; and the way vapour-compression, evaporative and fresh-air equipment interact.

Model predictive control (MPC) uses forecasts and a building model to trade energy cost
against comfort \cite{drgona2020,yao2021}. Deep reinforcement learning (DRL) learns a
policy from interaction with a simulator or a building \cite{yu2021,sierla2022}. Hybrid
schemes combine predictive and feedback layers \cite{arroyo2022}. Three gaps remain:
\begin{enumerate}
  \item Comparisons are often reported against weak baselines, or with results that
        cannot be reproduced. Many studies omit plant--model mismatch, forecast errors
        and repeated runs.
  \item Indian residential case studies that use typical-year weather, the national
        comfort range, ToD tariffs and the equipment actually installed in Indian homes
        are scarce.
  \item Thermal control and indoor air quality (IAQ) are rarely assessed together,
        although ventilation and filtration choices interact with cooling energy.
\end{enumerate}

This paper compares five controllers on the same validated simulation model of a
95\,m$^2$ Pune flat:
\begin{itemize}
  \item a split-AC thermostat;
  \item SIMC-tuned PID;
  \item economic MPC (E-MPC);
  \item a parameter-shared double DQN;
  \item a hybrid E-MPC/PI controller with a safety layer.
\end{itemize}
The contributions are:
\begin{itemize}
  \item An open, fully reproducible MATLAB/Octave benchmark of an Indian residence. It
        comprises a two-node-per-zone RC model discretised exactly; ISHRAE weather; an
        inverter VRF with part-load and latent behaviour; indirect evaporative cooling
        (IEC); a dedicated outdoor-air system (DOAS) with energy recovery; and moisture,
        CO$_2$ and PM$_{2.5}$ balances. A single script regenerates every number and
        figure.
  \item A like-for-like comparison under realistic information conditions:
        plant--model mismatch, sensor noise and forecast error. All controllers share
        one occupancy-based comfort schedule. Results are replicated over five seeds
        with confidence intervals and paired tests.
  \item A metric that separates cooling-induced discomfort from weather-driven cold,
        which no controller can remove because the dwelling has no heating, together
        with ToD cost, peak demand and peak-hour energy share.
  \item Robustness and ablation analyses: a +3\,K heatwave, thermal mass, tariff
        structure, model mismatch, MPC horizon, perfect information and DQN action
        resolution.
  \item An IAQ analysis of ventilation and filtration architectures, showing which
        design choices, independent of the thermal controller, govern CO$_2$, PM$_{2.5}$
        and humidity.
\end{itemize}
Section~\ref{sec:methods} describes the models and controllers,
Section~\ref{sec:results} presents the results, Section~\ref{sec:discussion} discusses
them, and Section~\ref{sec:limits} states the limitations.
